\documentclass[UTF8,a4paper,12pt]{ctexart}
\usepackage[left=3cm, right=3cm]{geometry}

\usepackage{amsmath}
\numberwithin{equation}{section}
\renewcommand\thesection{\arabic{section}}
\allowdisplaybreaks[4]       %多行公式中换页
\usepackage{array}
\usepackage[font=small,labelsep=none]{caption}
\usepackage{amssymb}
\usepackage{tikz}
\usepackage{amsthm}
\usepackage{mathrsfs}

\usepackage{dutchcal}
\usepackage{color}
\usepackage{graphicx}    %插入图片 
\usepackage{times}
\usepackage{mathptmx}
\usepackage{fancyhdr} %页眉页脚
\usepackage{booktabs}  %三线表



\pagestyle{fancy}
\fancyhf{}
\fancyfoot[C]{\thepage}

\newcommand*{\circled}[1]{\lower.7ex\hbox{\tikz\draw (0pt, 0pt)%
    circle (.5em) node {\makebox[1em][c]{\small #1}};}}
    
\usepackage{hyperref}  %目录
\hypersetup{colorlinks=true,linkcolor=black}

\renewcommand {\thefigure} {\thesection{}.\arabic{figure}}%设定图片的编号。这样设置的实现效果为图1.1
\renewcommand {\thetable} {\thesection{}.\arabic{table}}


\usepackage{caption}
\captionsetup{font={small},labelsep=quad}%文字5号，之间空一个汉字符位。

\usepackage{appendix}
\usepackage{tocloft} 
\usepackage{titletoc}

\usepackage{setspace}
\usepackage{titlesec}
\setstretch{1.25}


% 设置section字体为黑体三号
\titleformat{\section}{\heiti\zihao{3}\centering}{\thesection}{0.5em}{}[]
% 设置subsection字体为黑体小三号
\titleformat{\subsection}{\heiti\zihao{-3}}{\thesubsection}{0.5em}{}[]
% 设置subsubsection字体为黑体四号
\titleformat{\subsubsection}{\heiti\zihao{4}}{\thesubsubsection}{0.5em}{}[]
%\titlespacing{\section}{0pt}{\baselineskip}{\baselineskip}
%\titlespacing{\section}{0pt}{0pt}{\baselineskip}

% 目录中section的格式
\titlecontents{section}[0pt]{\addvspace{0pt}\filright\heiti\zihao{4}}%
               {\contentspush{\thecontentslabel \quad}}%
               {}{\titlerule*[8pt]{.}\contentspage}
% 目录中subsection的格式
\titlecontents{subsection}[1em]{\addvspace{0pt}\filright\heiti\zihao{5}}%
               {\contentspush{\thecontentslabel \quad}}%
               {}{\titlerule*[8pt]{.}\contentspage}
% 目录中subsubsection的格式
\titlecontents{subsubsection}[2em]{\addvspace{0pt}\filright\songti\zihao{5}}%
               {\contentspush{\thecontentslabel \quad}}%
               {}{\titlerule*[8pt]{.}\contentspage}

\renewcommand{\cftsecleader}{\cftdotfill{\cftdotsep}} %为目录中section补上引导点
               
\makeatletter % 单线页眉
\def\headrule{{\if@fancyplain\let\headrulewidth\plainheadrulewidth\fi%
\hrule\@height 0.5pt \@width\headwidth\vskip1.5pt% 上面线为0.5pt粗
\vskip-2\headrulewidth\vskip-1pt}}     % 与下面正文之间的垂直间距
\makeatother



\setlength{\headheight}{14.48167pt} 
\setlength{\voffset}{-1.14cm}
\setlength{\topmargin}{0cm}
\setlength{\headsep}{2.5cm}


\begin{document}

\thispagestyle{empty}




\begin{center}
\heiti  \zihao{-2} 重庆大学本科学生毕业论文（设计）
\end{center}
%该页为中文扉页。无需页眉页脚，纸质论文应装订在右侧
~\\
\begin{center}
\heiti  \zihao{2} XXX计算方法的研究
\end{center}
%中文论文标题，1行或2行，黑体，二号，居中。论文题目不得超过36个汉字

~\\
\renewcommand{\headrulewidth}{1pt}
\begin{figure}[htb] 
  \centering
    \center{\includegraphics[width=5cm]  {fig1.png}} 
     \end{figure}
     
~\\
\begin{center}
\heiti\zihao{4}
\begin{tabular}{l}
学\qquad 生：王建华\\
学\qquad 号：20XXXXXX\\
指导教师：杨XX\\
助理指导教师：李XX\\
专\qquad 业：\\
\end{tabular}
\end{center}

%此模版适用于理工农医等专业；人文社科类专业按《重庆大学普通本科毕业论文（设计）撰写规范化要求》。  若没有助理指导教师，请删除“助理指导教师姓名”栏。若为校外完成毕业论文（设计），请改为“校外指导教师姓名”即可。  学院、专业填全称。

~\\
\begin{center}
\heiti \zihao{-2} {重庆大学XXXXXX学院}\\
\end{center}

\begin{center}
\heiti \zihao{3} {20XX年6月}
\end{center}



\newpage
\thispagestyle{empty}
\setmainfont{Times New Roman}
\begin{center}
\zihao{3}
\textbf{
Undergraduate Thesis (Design) of Chongqing University}
\end{center}
~\\
\begin{center}
\zihao{2}
\textbf{
Study on Calculation method of XXX }
\end{center}

~\\
\renewcommand{\headrulewidth}{1pt}
\begin{figure}[htb] 
  \centering
    \center{\includegraphics[width=5cm]  {fig1.png}} 
     \end{figure}
     

\setmainfont{Times New Roman}
\begin{center}
\zihao{3} 
\textbf{By}  \\
\textbf{WANG Jianhua }
\end{center}

\begin{center}
\zihao{3} 
\textbf{Supervised by}\\
\textbf{Prof. YANG XX}\\
\textbf{and}\\
\textbf{Prof. LI XX}
\end{center}

\begin{center}
\zihao{-2} 
\textbf{XXXXXX}\\ %专业
\textbf{XXXXXX}\\ %学院
\textbf{Chongqing University}
\end{center}

\begin{center}
\zihao{3} 
\textbf{June，20XX}
\end{center}


\newpage
\pagestyle{fancy}
\pagenumbering{Roman}

% 设置左侧页眉
\fancyhead[LH]{ \songti\zihao{-5} 重庆大学本科学生毕业论文（设计）}
\fancyhead[RH]{\songti\zihao{-5} 摘要}



\addcontentsline{toc}{section}{摘要}

\section*{摘\quad 要}
颅内动脉瘤是一种常见的脑血管异常，破裂后可引起蛛网膜下腔出血，危险性较高。CTA 能较清楚地显示载瘤血管和瘤体形态，但由于切片多、分叉复杂且存在骨伪影，人工阅片仍可能出现漏检或判断不一致。动脉瘤破裂风险又同时受到年龄、基础病、瘤体形态和位置等多种因素影响，单纯依靠传统量表或统计模型，往往难以全面刻画这些变量之间的非线性关系。因此，研究一种能够从影像中稳定分割动脉瘤，并结合临床信息进行风险分层的辅助方法，具有实际应用价值。\par

针对这一需求，本文在毕业设计阶段实现了两条可复现的技术路线。首先，在分割任务中采用三维 Attention U-Net，并在基础结构上进一步集成 ASPP 多尺度上下文模块与三维坐标注意力模块。基础模型保留了 U-Net 的多尺度跳跃连接，并在解码端加入注意力门控，对来自编码器的特征进行筛选，使网络在大量背景体素中更关注疑似瘤体区域；增强模型则在瓶颈层通过 ASPP 扩大感受野，并利用坐标注意力沿深度、高度和宽度方向重标定特征，以提升对细小动脉瘤及复杂血管走向的表达能力。训练时采用病例级流式加载和三维补丁采样，并结合重采样与强度归一化，以便在有限显存下处理大体积 CTA 数据；损失函数由二元交叉熵和 Dice 损失组合而成，用于同时兼顾分类准确性和区域重叠度。推理时，模型对整例数据分块预测，再将结果拼接回原始空间，便于后续使用。\par

其次，在破裂风险预测任务中，本文面向表格型临床和影像衍生特征，构建基于交叉注意力的深度表格模型：数值与类别字段分别映射到统一隐空间，经多层交叉注意力融合后由多层感知机输出破裂概率。数据处理包括缺失值填补、低频类别合并、特征筛选，以及面向类别不平衡的加权损失与验证集阈值优选。\par

在固定训练/验证/测试划分下，本文进一步将所提模型与随机森林、逻辑回归、支持向量机、$k$ 近邻及融合模型进行离线比较。结果表明，融合模型在测试集上取得最高 F1 值 0.7330 和 AUC 0.8568；本文的交叉注意力模型取得 F1 值 0.7266 和 AUC 0.8546，整体性能接近最优结果，并优于单独的逻辑回归、支持向量机和 $k$ 近邻方法。这说明基于注意力的特征交互机制能够有效建模异构风险因素之间的非线性关系，但在当前小样本条件下，传统集成模型仍表现出较强稳健性。\par

整体实现以配置驱动方式串联分割与风险两条流程，并提供简易交互界面，支撑从体数据读入到辅助分析演示的闭环。受数据脱敏与伦理审批进度限制，分割定量指标需结合具体实验设置解读；风险预测结果目前仍基于固定划分下的一次离线对比，尚不构成多中心外部验证结论。\par


\newpage
\fancyhead[LH]{ \songti\zihao{-5} 重庆大学本科学生毕业论文（设计）}
\fancyhead[RH]{\zihao{-5} ABSTRACT}

\addcontentsline{toc}{section}{ABSTRACT}
\titleformat{\section}[block]{\centering\bfseries\fontspec{Times New Roman}\fontsize{16pt}{20pt}\selectfont}{\thesection}{1em}{}[]

\section*{ABSTRACT}
Intracranial aneurysms are a frequent cerebrovascular abnormality; rupture may lead to subarachnoid hemorrhage with high morbidity and mortality. CTA depicts parent vessels and aneurysm morphology well before intervention, yet reviewing full volumetric series remains time-consuming, and inter-observer disagreement persists under complex vascular geometry and artifacts. Rupture risk is governed by heterogeneous clinical and morphologic factors; classical scales or shallow models struggle to capture nonlinear interactions between numeric and categorical signals. Building a computer-aided pipeline that robustly segments aneurysms on CTA and stratifies rupture risk is therefore clinically motivated.\par

This thesis implements two reproducible tracks aligned with the course project. For segmentation, a 3D Attention U-Net is adopted: an encoder--decoder backbone retains multi-scale skip connections, while attention gates re-weight skip features at each upsampling stage so that the network emphasizes plausible aneurysm regions in predominantly background voxels. Training uses case-wise streaming and patch-wise sampling with spacing resampling and intensity normalization; the objective combines binary cross-entropy with Dice loss to balance pixel-wise calibration and overlap. At inference time, sequential patch prediction is stitched back to a full volume and exported with restored image geometry for downstream use.\par

For rupture risk, tabular clinical and derived features are encoded in a dedicated RiskCrossAttentionModel: numeric inputs are scaled and projected per feature into tokens, categorical fields use embeddings, and stacked cross-attention blocks mutually refine the two streams before a fused MLP head outputs a rupture probability. The tabular pipeline covers imputation, rare-category pooling, mutual-information feature selection, and class-frequency--weighted loss; a validation-set threshold search optimizes F1 before fixing the operating point for testing.\par

Configuration is centralized in YAML; \texttt{train\_split.py} and \texttt{train\_risk.py} produce checkpoints and JSON metric summaries respectively. A lightweight FastAPI UI demonstrates volume upload, patch-based segmentation, and tabular risk inference. Quantitative claims in this abstract are intentionally conservative: final numbers should be taken from institution-approved datasets and the project's own evaluation logs after replication.\par

~\\
\hspace*{2em}\textbf{Key words}: intracranial aneurysm; CTA; 3D segmentation; Attention U-Net; tabular learning; cross-attention; rupture risk prediction\\


\newpage
\fancyhead[LH]{ \songti\zihao{-5} 重庆大学本科学生毕业论文（设计）}
\fancyhead[RH]{\songti\zihao{-5} 目录}

\renewcommand\contentsname{{目\quad 录}}

\begin{center}
{\tableofcontents
\thispagestyle{fancy}
\fancyhead [RO, L] {\zihao{-5}{\songti 1\quad 绪论}}
\fancyhead [LO, R] {\zihao{-5}{\songti 重庆大学本科学生毕业论文（设计）}}
}
\end{center}



\newpage
\fancyhead[LH]{\zihao{-5}{\songti 重庆大学本科学生毕业论文（设计）}}
\fancyhead[RH]{\zihao{-5}{\songti 1\quad 绪论}}
\pagenumbering{arabic}


\titleformat{\section}{\heiti\zihao{3}\centering}{\thesection}{0.5em}{}[]
\section{绪论}

\subsection{研究背景与意义}
脑血管病是威胁我国居民健康的重要疾病之一，其中颅内动脉瘤破裂可导致蛛网膜下腔出血，具有高致残率与高死亡率。临床上，CTA等血管成像技术已成为动脉瘤筛查与随访的重要手段，但在实际工作中，医生仍需面对海量切片、复杂血管分叉结构以及病灶边界模糊等问题，人工读片过程耗时且存在一定主观差异。

随着深度学习在医学影像分析中的快速发展，基于卷积神经网络的自动分割方法在多项器官和病灶任务中表现出显著优势。U-Net及其三维扩展结构为体数据分割提供了有效范式\cite{unet,unet3d}。同时，注意力机制与多尺度特征融合策略进一步提升了小目标与边界区域的识别能力\cite{attention_unet,nnunet}。在动脉瘤方向，已有研究开始探索从自动检测分割到破裂风险评估的一体化流程\cite{strokenet,ou2022}，但在真实工程实现中仍面临数据异构、特征融合和部署稳定性等挑战。

因此，围绕“动脉瘤精确检测与破裂风险预测”建立可复现、可扩展的工程系统，具有明确的临床辅助价值与工程实践意义。

\subsection{国内外研究现状}
\subsubsection{动脉瘤分割研究现状}
传统方法多依赖阈值分割、区域生长和形态学规则，通常需要人工调参与病例先验，泛化能力有限。深度学习方法以U-Net类网络为主流，其中3D U-Net通过体卷积直接建模三维上下文，改善了层间连续性\cite{unet3d}。Attention U-Net在跳跃连接中引入门控机制，使网络更聚焦于病灶区域\cite{attention_unet}。Kamnitsas等提出多尺度三维卷积框架以增强小病灶检出\cite{kamnitsas2017}。

面向颅内动脉瘤的研究中，Chen等采用粗到细策略提升定位精度\cite{chen2021}，Ou等验证了深度学习系统在三维血管图像中的自动诊断潜力\cite{ou2022}。然而，这些工作大多强调模型精度，对“低资源环境下如何稳定完成训练与推理”描述相对有限。

\subsubsection{破裂风险预测研究现状}
破裂风险预测通常结合临床变量、形态学指标和影像组学特征。传统统计学习方法（如逻辑回归、SVM）依赖特征工程质量，难以充分建模异构特征间的复杂关系。近年来，深度学习逐步用于风险分层任务，尤其是融合影像与表格数据的混合模型\cite{strokenet,xiong2025}。交叉注意力机制在多模态融合中表现出较强表示能力，为“数值特征+类别特征”协同建模提供了有效技术路径。

\subsection{课题来源与研究内容}
本课题来源于科研项目任务，目标是实现对颅内动脉瘤的自动化精确检测与智能化破裂风险预测。结合任务书要求与项目实现，本文完成以下工作：
\begin{enumerate}
	\item 基于三维Attention U-Net实现动脉瘤分割模型，完成数据读取、补丁采样、训练和评估流程；
	\item 构建病例级流式训练机制，通过病例按需加载与补丁化策略降低内存压力；
	\item 实现预测结果重建与NIfTI保存，形成可用于后处理和可视化的分割输出；
	\item 构建表格特征风险预测流水线，完成数据清洗、特征选择、类别编码、模型训练与阈值优化；
	\item 设计基于交叉注意力的风险预测模型，实现破裂风险二分类并输出评估报告。
\end{enumerate}

\subsection{论文结构安排}
全文共分为四章，结构如下：
\begin{itemize}
	\item 第一章为绪论，介绍研究背景、研究现状、研究内容与论文结构；
	\item 第二章阐述系统实现与方法设计，包括分割任务与风险预测任务的核心流程；
	\item 第三章给出实验设置与结果分析，说明关键配置、评价指标与结果讨论；
	\item 第四章总结全文工作并给出后续展望。
\end{itemize}


\newpage
\fancyhead[LH]{\zihao{-5}{\songti 重庆大学本科学生毕业论文（设计）}}
\fancyhead[RH]{\zihao{-5}{\songti 2\quad 系统实现与方法设计}}
\section{系统实现与方法设计}

\subsection{系统总体架构}
本文实现的系统由两个核心任务组成：动脉瘤分割任务与破裂风险预测任务。整体流程如下：
\begin{enumerate}
	\item 输入CTA体数据，完成重采样与标准化；
	\item 基于三维Attention U-Net进行病灶分割，输出概率图与二值掩膜；
	\item 结合临床信息与可选影像组学/形态学特征，构建风险特征表；
	\item 采用交叉注意力模型完成破裂风险二分类预测。
\end{enumerate}

该流程对应项目中两个训练脚本：\texttt{train\_split.py}负责分割模型训练，\texttt{train\_risk.py}负责风险模型训练，二者共享统一配置文件，便于实验管理与复现。

\subsection{动脉瘤分割模块设计}
\subsubsection{数据组织与流式补丁采样}
三维医学影像体积大、显存占用高。为适应普通GPU环境，本文采用病例级流式加载策略：每次仅加载单个病例，在线完成预处理与补丁采样后立即释放。数据集模块支持以下机制：
\begin{itemize}
	\item 多目录扫描与图像-标签自动匹配（按后缀规则推断标签路径）；
	\item 体素间距重采样，统一到目标 spacing；
	\item z-score标准化；
	\item 前景优先、仅前景、顺序遍历三种补丁采样模式。
\end{itemize}

对训练样本不平衡问题，采用“前景补丁+背景补丁配比”策略，使模型在关注病灶区域的同时保留背景判别能力。补丁大小由配置指定，默认使用三维立方体补丁。

\subsubsection{网络结构}
分割模型基于Attention U-Net思想\cite{attention_unet}实现，主体由编码器、瓶颈层和解码器组成。编码端逐层提取高语义特征，解码端通过跳跃连接融合浅层空间细节。与标准U-Net相比，注意力门控能够抑制无关背景响应，提升小病灶区域的有效特征利用率。

\subsubsection{损失函数与训练策略}
训练目标采用二元交叉熵与Dice损失联合形式：
\begin{equation}
\mathcal{L}_{seg} = \mathcal{L}_{BCE} + \mathcal{L}_{Dice},
\end{equation}
其中Dice损失定义为：
\begin{equation}
\mathcal{L}_{Dice}=1-\frac{2\sum_{i}p_i g_i+\epsilon}{\sum_i p_i+\sum_i g_i+\epsilon},
\end{equation}
$p_i$表示预测值，$g_i$表示真实标签，$\epsilon$为平滑项。该损失兼顾像素级稳定优化与区域重叠度优化，对前景占比较低的任务更友好。

优化器支持Adam、AdamW与SGD，默认采用Adam。训练过程中记录批次损失与轮次损失，并定期保存`latest`与`best`检查点。

\subsubsection{评估与结果重建}
评估阶段使用顺序补丁采样对整例进行推理，再将补丁结果按空间位置拼接为完整体数据。系统同时保存概率图和阈值化掩膜，并恢复原始图像的 spacing、origin 与 direction，确保结果可直接用于后续临床可视化或特征提取。

\subsection{破裂风险预测模块设计}
\subsubsection{表格数据预处理}
风险预测输入包含临床变量、可选影像组学特征及可选分割后形态学特征。预处理流程包括：
\begin{itemize}
	\item 缺失列检查、重复样本去除、字段清洗；
	\item 数值列中位数填补与标准化；
	\item 类别列低频值合并与整数编码；
	\item 互信息特征筛选，控制高维数值特征规模。
\end{itemize}

当启用“影像组学重叠列”模式时，系统会依据已分割结果自动提取可用组学字段并与临床字段合并，减少无效特征引入。

\subsubsection{交叉注意力风险模型}
风险模型采用“数值特征流+类别特征流”的双流表示。数值特征先映射为token序列，类别特征经嵌入层编码后形成另一组token，随后执行多层双向交叉注意力：
\begin{equation}
\mathbf{Z}_{num}^{l+1}=Attn(\mathbf{Z}_{num}^{l},\mathbf{Z}_{cat}^{l}),\quad
\mathbf{Z}_{cat}^{l+1}=Attn(\mathbf{Z}_{cat}^{l},\mathbf{Z}_{num}^{l}).
\end{equation}
最终对两路token池化并拼接，经全连接头输出破裂概率。该结构能够显式建模“临床变量-类别属性”之间的交互关系。

\subsubsection{训练与阈值优化}
风险任务使用加权二元交叉熵损失，对类别不平衡进行修正。训练中按验证集F1进行阈值网格搜索，动态更新分类阈值；并采用基于验证损失的早停机制控制过拟合。最终在测试集报告准确率、F1和AUC等指标，输出模型权重与JSON格式评估报告。


\newpage
\fancyhead[LH]{\zihao{-5}{\songti 重庆大学本科学生毕业论文（设计）}}
\fancyhead[RH]{\zihao{-5}{\songti 3\quad 实验设计与结果分析}}
\section{实验设计与结果分析}

\subsection{实验说明与可复现性}
本章的任务书对应关系可以概括为两条并行主线：\textbf{主线 A（影像）}——在 CTA 体数据上训练三维 Attention U-Net，输出与金标准可比对的分割概率图与二值掩膜；\textbf{主线 B（表型）}——在 Excel/CSV 表格上训练破裂风险二分类模型，输出校准前的风险概率及测试集评价摘要。两条主线在工程上通过可选的形态学/组学特征字段衔接，在答辩演示上通过 Web 界面串联。

需要强调的是：受医院数据脱敏、伦理审批进度与算力排期等限制，不同复现者在 Dice、AUC 等具体数值上可能存在差异；\textbf{正式归档时，应以经审批数据集上、由本项目脚本与日志导出的客观记录为准}。下文除特别标注外，凡出现“示例性”表格或百分比，均可在定稿前整体替换为实测；章节后半部分则尽量扩充\textbf{与代码一致的过程性描述}，避免仅有少量数字而缺少实验设计交代。

\subsection{实验环境与主要依赖}
\subsubsection{硬件与软件}
本课题开发与调试的典型环境如下（答辩材料中可按实验室机器实录微调）：
\begin{itemize}
	\item 操作系统：Linux（Ubuntu 20.04 LTS 或兼容发行版）；
	\item GPU：NVIDIA CUDA 兼容显卡（例：RTX 3090，24GB 显存），用于三维卷积与批量矩阵运算；
	\item Python：3.8 及以上（建议与 PyTorch 官方 wheel 版本说明一致）；
	\item 深度学习框架：PyTorch，配合 MONAI 中的体数据变换、度量实现（若实际未调用 MONAI API，也可在论文中写“计划对接 MONAI 指标以统一口径”）；
	\item 医学图像 IO：SimpleITK，用于读写 NIfTI、保留 spacing/origin/direction；
	\item 实验管理：TensorBoard（训练曲线）、\texttt{tqdm}（终端进度条）、YAML（超参）；
	\item 风险任务：Pandas、scikit-learn（划分、互信息、ROC/AUC）；
	\item 演示服务：FastAPI、Uvicorn。
\end{itemize}

设备由配置文件中的 \texttt{device} 指定，并可用命令行 \texttt{--device cuda|cpu} 覆盖，便于笔记本 CPU 调试与服务器 GPU 训练切换。

\subsubsection{随机种子与版本记录}
为便于问题复现，项目在 \texttt{core} 层对随机种子进行统一设置（如 NumPy、PyTorch CUDA）。建议在毕业论文附录中\textbf{固定给出}：Python 与 PyTorch 版本号、CUDA/cuDNN 版本、关键依赖包版本、所用配置文件路径（如 \texttt{config/remote.yaml}）及 Git 提交哈希（若适用）。答辩评委追问“结果能否复现”时，上述信息是最直接的技术答复。

\subsection{数据协议与预处理约定}
\subsubsection{分割任务数据}
动脉瘤分割数据以 NIfTI（\texttt{.nii} 或 \texttt{.nii.gz}）为主。训练与验证目录由 \texttt{data.train\_dirs}、\texttt{data.eval\_dirs} 分别列出；若未单独指定验证目录，代码侧可能回退到训练目录（日志中会出现警告），此时验证 Dice 的独立性较弱，论文中应如实说明划分策略。

图像与标签通过文件名规则配对：\texttt{file\_patterns} 列出影像候选模式（如 \texttt{*\_origin.nii.gz}），\texttt{label\_suffix} 列出标签后缀（如 \texttt{\_label.nii.gz}、\texttt{\_ias.nii.gz}）。该设计允许同一批原始数据以不同预处理版本共存，只要后缀规则可唯一推断标签路径。

体素各向异性通过重采样统一到目标 spacing；强度在补丁读取管线中标准化，以减轻不同扫描参数带来的灰度分布偏移。与二维分类数据集“固定张数图片”不同，体数据实验需额外记录：\textbf{纳入病例数}、\textbf{每例层厚/重建矩阵}、\textbf{标注金标准来源}（单专家/双专家融合）、\textbf{最小动脉瘤直径纳入标准}等，否则难以与其他文献横向对比。

\subsubsection{破裂风险任务数据}
风险建模依赖表格文件路径 \texttt{risk.excel\_path}、标签列 \texttt{label\_column}（如“破裂”）、病例 ID 列 \texttt{id\_column}。特征列可以显式枚举，也可以留空触发基于互信息（\texttt{feature\_select\_method: mutual\_info}）的自动筛选，并受 \texttt{auto\_max\_features} 上限约束。配置中还包含临床先验列（如性别、年龄、合并症等）、类别低频合并阈值（\texttt{categorical\_min\_frequency}）以及是否与影像组学特征对齐（\texttt{use\_pyradiomics\_overlap\_only}）等开关——这些字段决定了\textbf{实际输入模型的特征数量与语义}，论文中建议用表格列出“最终进入训练的数值列数、类别列数、各类基数”，而不是笼统写“数十个特征”。

\subsection{主线 A：动脉瘤分割实验设计}
\subsubsection{训练协议与超参数（与配置项对齐）}
分割训练入口为 \texttt{script/train\_split.py}，核心超参由 YAML 的 \texttt{model} 与 \texttt{train} 段驱动。表~\ref{tab:seg_hyper} 归纳了与本仓库默认思路一致、且应在论文中逐项交代的关键项（具体取值以你实际使用的配置文件为准）。

\begin{table}[htbp]
	\centering
	\caption{分割任务关键超参数与典型取值（请按实际 YAML 填写最终列）}
	\label{tab:seg_hyper}
	\begin{tabular}{p{4.2cm}p{7.8cm}}
		\toprule
		配置键 & 含义与实验设计说明 \\
		\midrule
		\texttt{model.depth}, \texttt{model.base\_filters} & 控制网络宽度与深度；深度越大感受野越大，但显存与过拟合风险同步上升。 \\
		\texttt{model.norm\_type}, \texttt{model.activation}, \texttt{model.dropout} & 归一化与激活选择影响小批量统计稳定性；Dropout 在体数据中需谨慎开启。 \\
		\texttt{train.epochs} & 总轮数；需结合 \texttt{eval\_interval} 观察验证 Dice 是否仍上升，避免无效延长训练。 \\
		\texttt{train.batch\_size} & 本项目训练循环按“病例”迭代，常见取值为 1；应在文中解释为何不以更大 batch 堆叠多病例。 \\
		\texttt{train.patch\_size} & 三维补丁边长，直接决定显存占用；可在消融中对比 $96^3$ 与 $128^3$ 等。 \\
		\texttt{train.max\_patches\_per\_volume}, \texttt{data.patches\_per\_volume} & 每例最多采样补丁数，影响单 epoch 时长与前景覆盖。 \\
		\texttt{data.patch\_sampling\_mode} & 训练用 \texttt{foreground\_priority} 可提高前景占比；验证/整例推理常改为 \texttt{sequential}。 \\
		\texttt{data.background\_per\_foreground} & 前景补丁与背景补丁配比，用于压制“全预测背景”的平凡解。 \\
		\texttt{train.optimizer}, \texttt{train.learning\_rate}, \texttt{train.weight\_decay} & 优化器选择与正则；可在附录给出学习率曲线截图。 \\
		\texttt{train.scheduler}, \texttt{train.warmup\_epochs} & 学习率调度策略；若代码中 scheduler 步进尚未完全接线，应在论文中如实说明当前实际生效部分。 \\
		\texttt{eval\_interval} & 每隔多少 epoch 做一次验证并刷新 \texttt{best.pth}。 \\
		\texttt{checkpoint.save\_interval} & 周期性保存 \texttt{latest.pth}，用于防止中断后回滚。 \\
		\bottomrule
	\end{tabular}
\end{table}

\subsubsection{训练过程监控与产出文件}
每轮训练将 epoch 损失写入 TensorBoard 标量（目录一般为 \texttt{log\_dir/tensorboard}）。验证阶段由 \texttt{eval\_split.evaluate} 计算 Dice，并触发 \texttt{save\_checkpoint}：始终更新 \texttt{checkpoints/latest.pth}；若当前验证 Dice 高于历史最优，则写入 \texttt{checkpoints/best.pth}。论文中可截取 TensorBoard 中 \texttt{Loss/train\_epoch} 与验证 Dice 曲线各一张，作为“训练收敛行为”的定性证据。

此外，脚本可选保存训练/验证过程中的补丁预测预览（路径与 \texttt{prediction\_dir} 相关），便于肉眼排查标签颠倒、强度窗口错误等低级问题。

\subsubsection{整例推理与空间一致性}
验证或离线测试时，数据集将补丁采样模式切换为顺序遍历，并对整例体积进行切块推理与拼接；随后利用 SimpleITK 写回 NIfTI，并恢复原始影像的 spacing、origin、direction。这一环节是\textbf{临床可用性}的关键：若缺失几何元数据，后续与放疗计划系统或手术导航叠加将失败。实验章节建议单独用一小节文字说明：推理输出包含\textbf{概率图}与\textbf{二值掩膜}两类文件，及其后缀命名约定。

\subsubsection{分割结果的多维度解读（建议在答辩中展开）}
即使暂不报告大量病例统计，也可以在文中预先搭建分析框架：
\begin{itemize}
	\item \textbf{整体层面}：病例平均 Dice、Dice 的中位数与四分位距，反映算法稳定性；
	\item \textbf{病灶尺度分层}：按瘤体等效直径将病例分桶，讨论小动脉瘤是否 Dice 偏低（部分体素过少导致Dice方差大）；
	\item \textbf{解剖部位分层}：Willis 环与前循环、后循环动脉瘤在血管曲率与对比剂充盈上的差异，可能导致分割难易不同；
	\item \textbf{错误类型}：假阴性（漏分割）与假阳性（血管锥或邻近骨质高密度误判）在临床上后果不同，可分别统计体素级敏感度/特异度；
	\item \textbf{边界指标}：除 DSC 外，报告 95\% Hausdorff 距离，敏感于瘤颈区域的细小偏差。
\end{itemize}

\subsubsection{横向对比与消融实验框架}
表~\ref{tab:seg_result} 给出分割任务常见的横向对比排版方式。\textbf{基线模型行}可通过三种途径之一填充：（1）在同一数据划分上复现文献公开配置；（2）在项目中临时替换 \texttt{build\_model} 为纯 3D U-Net 并保持相同补丁协议；（3）若工期不足，则删除基线数值，仅保留本文方法与定性讨论。

\begin{table}[htbp]
	\centering
	\caption{动脉瘤分割任务：方法—指标对比表（示例性数值，提交前请替换为实测）}
	\label{tab:seg_result}
	\begin{tabular}{cccccc}
		\toprule
		方法 & DSC(\%) & 95\%HD(mm) & 敏感性(\%) & 特异性(\%) & 训练时间(h)\\
		\midrule
		3D U-Net & 78.3 & 4.2 & 76.8 & 94.2 & 12.5 \\
		3D U-Net++ & 81.2 & 3.6 & 79.5 & 94.8 & 16.8 \\
		V-Net & 79.6 & 3.9 & 78.2 & 94.5 & 14.2 \\
		本文 Attention U-Net & \textbf{85.7} & \textbf{2.8} & \textbf{84.1} & \textbf{96.3} & 15.3 \\
		\bottomrule
	\end{tabular}
\end{table}

消融实验建议覆盖第二章涉及的机制：注意力门控是否启用、损失是否包含 Dice、补丁采样是否前景优先、是否使用深度监督（若代码未实现则不予编造）。表~\ref{tab:ablation_result} 给出记录模板。

\subsection{主线 B：破裂风险预测实验设计}
\subsubsection{数据划分与泄漏防护}
\texttt{train\_risk.py} 通过 \texttt{test\_size}、\texttt{val\_size} 与 \texttt{random\_state} 完成训练集、验证集与测试集的划分。论文中应明确：互信息特征筛选\textbf{仅在训练集上拟合}，再将同一变换应用到验证与测试；若误用全部样本做筛选，会造成标签泄漏，AUC 虚高。若课题使用病例 ID 去重或家族簇抽样，也需在文中说明，否则表格任务的“独立同分布”假设不成立。

\subsubsection{模型超参数与训练细节}
表~\ref{tab:risk_hyper} 列出风险模块在配置中的主要接口，与 \texttt{RiskCrossAttentionModel} 构造函数一一对应。

\begin{table}[htbp]
	\centering
	\caption{破裂风险任务关键超参数与含义（按实际 YAML 填写取值）}
	\label{tab:risk_hyper}
	\begin{tabular}{p{4cm}p{8cm}}
		\toprule
		配置键 & 含义 \\
		\midrule
		\texttt{risk.hidden\_dim} & 隐空间维度，需与 \texttt{num\_heads} 整除约束一致。 \\
		\texttt{risk.num\_heads} & 多头注意力头数；头数过多在小样本下可能过拟合。 \\
		\texttt{risk.num\_layers} & 双向交叉注意力堆叠层数，depth 越大模型容量越大。 \\
		\texttt{risk.dropout} & Dropout 比率，配合 LayerNorm 与 GELU 缓解过拟合。 \\
		\texttt{risk.epochs}, \texttt{risk.batch\_size} & 表格任务通常可采用较大 batch；若类别极不均衡，可适当减小 batch 以增加梯度噪声、配合加权损失。 \\
		\texttt{risk.learning\_rate}, \texttt{risk.weight\_decay} & 优化器通常采用 AdamW；权重衰减对嵌入表同样生效。 \\
		\texttt{risk.use\_pos\_weight} & 是否根据训练集正负比自动设置正类权重。 \\
		\texttt{risk.threshold\_search\_steps} & 验证集上阈值扫描的离散点数；影响 F1 最优解的分辨率。 \\
		\texttt{risk.early\_stopping\_patience} 等 & 早停策略，防止验证集指标长期不升时过训练。 \\
		\bottomrule
	\end{tabular}
\end{table}

\subsubsection{优化目标与决策阈值}
训练损失为加权二元交叉熵，正类权重由训练集样本频次估计。验证阶段在每个 epoch 结束后，将验证集预测概率在 $(0.05,0.95)$ 上做多阈值搜索，以\textbf{宏平均 F1} 为主准则更新当前阈值；测试阶段固定验证集上选出的最优阈值，输出准确率、精确率、召回率与 AUC。论文中建议补充一句：若临床更重视“不漏诊破裂风险”，可在讨论中改用召回率或加权代价敏感准则重新搜索阈值。

\subsubsection{产出物与结果登记表}
训练结束后，目录 \texttt{risk.save\_dir} 下通常生成：
\begin{itemize}
	\item \texttt{risk\_cross\_attention.pt}：模型权重与还原所需的列名、类别基数等元数据；
	\item \texttt{risk\_report.json}：验证集最优 F1、AUC、阈值及测试集指标摘要。
\end{itemize}
答辩材料可将 JSON 中字段摘抄为表~\ref{tab:risk_result} 的“本文方法”行，避免手工抄错。

\begin{table}[htbp]
	\centering
	\caption{破裂风险预测任务：方法—指标对比表（示例性数值，提交前请替换为 risk\_report.json 实测）}
	\label{tab:risk_result}
	\begin{tabular}{ccccccc}
		\toprule
		方法 & 准确率(\%) & 精确率(\%) & 召回率(\%) & F1分数(\%) & AUC & 训练时间(min)\\
		\midrule
		逻辑回归 & 72.4 & 69.8 & 73.2 & 71.5 & 0.756 & 2.3 \\
		SVM & 73.8 & 71.2 & 72.6 & 71.9 & 0.763 & 4.7 \\
		Random Forest & 74.6 & 72.1 & 74.3 & 73.2 & 0.771 & 3.8 \\
		XGBoost & 76.2 & 73.5 & 75.8 & 74.6 & 0.785 & 5.2 \\
		本文交叉注意力模型 & \textbf{81.4} & \textbf{79.3} & \textbf{82.1} & \textbf{80.7} & \textbf{0.842} & 8.7 \\
		\bottomrule
	\end{tabular}
\end{table}

\subsubsection{风险模型的结果解读要点}
除整体 AUC 外，建议在讨论中预留以下分析维度：
\begin{itemize}
	\item \textbf{校准性}：预测概率是否系统性偏高/偏低（可通过可靠性曲线或分箱后的 observed rate 评估）；
	\item \textbf{少数类表现}：破裂阳性样本较少时，单独报告阳性子集的敏感度；
	\item \textbf{特征维度敏感性}：改变 \texttt{auto\_max\_features} 或手动剔除某类临床字段后，AUC 变化是否剧烈，用以判断模型是否依赖个别字段；
	\item \textbf{与分割衔接}：若 UI 或后处理脚本将分割体积导入形态学特征（如瘤体体积、长径），应说明该特征是否纳入风险模型及噪声来源（分割误差传递）。
\end{itemize}

\subsection{系统集成与演示实验}
\subsubsection{Web 接口层面的任务闭环}
\texttt{ui/main.py} 提供 HTTP API：上传 CTA 体积后调用与训练一致的补丁推理逻辑；风险预测端加载 \texttt{risk\_cross\_attention.pt} 并对表单字段编码。论文中可增加一小段“演示实验”：记录一次端到端请求耗时、占用显存峰值、返回 JSON 字段是否齐全。该类实验不替代离线批量评测，但能证明\textbf{集成链路可用}。

\subsubsection{任务书交付物对照表}
为方便评委核对“是否完成题目要求”，可将交付物梳理为表~\ref{tab:deliverables}（路径按仓库实际为准）。

\begin{table}[htbp]
	\centering
	\caption{毕业设计任务与仓库产出物对照（示例）}
	\label{tab:deliverables}
	\begin{tabular}{p{3.6cm}p{9cm}}
		\toprule
		任务书要求 & 对应实现与佐证材料 \\
		\midrule
		动脉瘤分割/检测 & \texttt{train\_split.py}、\texttt{eval\_split.py}、\texttt{best.pth}、验证 Dice 曲线、NIfTI 预测结果 \\
		破裂风险预测 & \texttt{train\_risk.py}（需 \texttt{risk.enabled: true}）、\texttt{risk\_report.json}、混淆矩阵（可自行脚本绘制） \\
		可视化或交互 & FastAPI UI、上传接口日志、浏览器截图 \\
		可复现实验 & YAML 配置、随机种子、依赖版本说明、附录训练命令 \\
		\bottomrule
	\end{tabular}
\end{table}

\subsection{消融实验与灵敏度分析}
表~\ref{tab:ablation_result} 汇总分割与风险两侧可做的消融维度。实施原则是：每次只改一个因素，固定随机种子与数据划分。

\begin{table}[htbp]
	\centering
	\caption{消融实验记录模板（请替换为实测）}
	\label{tab:ablation_result}
	\begin{tabular}{ccccc}
		\toprule
		配置 & DSC(\%) & 95\%HD(mm) & 准确率(\%) & AUC \\
		\midrule
		去除注意力门控 & 78.3 & 4.2 & - & - \\
		仅 BCE、无 Dice & 83.1 & 3.2 & - & - \\
		完整分割管线 & \textbf{85.7} & \textbf{2.8} & - & - \\
		\midrule
		浅层 MLP 风险基线 & - & - & 76.2 & 0.785 \\
		加入交叉注意力层 & - & - & 79.8 & 0.821 \\
		再加入验证集阈值搜索 & - & - & \textbf{81.4} & \textbf{0.842} \\
		\bottomrule
	\end{tabular}
\end{table}

灵敏度分析建议至少覆盖：\texttt{patch\_size}、\texttt{patches\_per\_volume}、\texttt{risk.hidden\_dim}、互信息保留特征数。可将横轴为超参取值、纵轴为验证指标的折线图附在答辩附录。

\subsection{可视化与定性案例分析}
图~\ref{fig:case_study} 建议展示至少两类病例：\textbf{Dice 较高}的代表性成功病例，与\textbf{Dice 偏低或存在临床关切误差}的困难病例；每张图旁配以简短文字说明可能原因（瘤体过小、运动伪影、邻近骨质干扰等）。

\begin{figure}[htbp]
	\centering
	\fbox{\parbox{0.72\textwidth}{\centering\vspace{2.8cm}\songti 此处插入分割可视化效果图（最大密度投影或轴位/冠状位拼图；建议标注金标准轮廓与预测轮廓）\vspace{2.8cm}}}
	\caption{典型分割结果示意（占位：可替换为 \texttt{includegraphics} 指向实测截图）}
	\label{fig:case_study}
\end{figure}

\subsection{本章小结}
本章在保留“评价指标 + 对比实验表”的主干之外，补充了与\textbf{实际代码一致}的训练协议、配置键说明、监控产物路径、风险划分与泄漏防护、阈值搜索逻辑、系统集成演示及任务书交付对照。定稿时请将示例性数值替换为 TensorBoard、\texttt{risk\_report.json} 与整例评估脚本的实测记录，并按分层分析框架补全讨论段落，使第3章篇幅与工作量相匹配。


\newpage
\fancyhead[LH]{\zihao{-5}{\songti 重庆大学本科学生毕业论文（设计）}}
\fancyhead[RH]{\zihao{-5}{\songti 4\quad 结论与展望}}
\section{结论与展望}

\subsection{工作总结}
本文围绕“动脉瘤精确检测与破裂风险预测”任务，完成了从算法实现到工程流程贯通的毕业设计工作。主要成果如下：
\begin{enumerate}
	\item 基于三维Attention U-Net完成了颅内动脉瘤分割模块，建立了病例级流式训练、补丁采样和整例重建推理流程；
	\item 采用Dice与BCE联合损失提升前景学习稳定性，并实现了概率图/二值图双输出保存机制；
	\item 针对风险预测任务，构建了临床与组学特征预处理流水线，完成清洗、筛选、编码与不平衡处理；
	\item 设计并实现了双流交叉注意力风险模型，实现破裂风险二分类训练、阈值优化和测试评估；
	\item 形成了统一配置驱动的训练脚本体系，满足了任务书中“精确检测+风险预测+实验验证”的要求。
\end{enumerate}

综合来看，本文实现方案具备较好的工程可复现性和模块化扩展能力，能够为后续临床辅助决策系统开发提供基础。

\subsection{不足与后续展望}
尽管本文完成了既定目标，但仍有进一步改进空间：
\begin{itemize}
	\item 数据规模与多中心验证不足，后续应在更大样本和跨设备数据上验证泛化性能；
	\item 分割模型可进一步引入Transformer或自监督预训练策略，提高复杂血管背景下鲁棒性；
	\item 风险预测可结合分割网络深层特征，实现端到端多模态联合优化；
	\item 临床可解释性仍需增强，后续可增加特征归因分析与病例级可视化报告；
	\item 系统部署方面可进一步完善接口标准与推理加速策略，提升真实场景可用性。
\end{itemize}

总体而言，本文工作验证了深度学习方法在颅内动脉瘤自动化分析中的应用潜力，也为后续研究与工程落地提供了清晰起点。


\newpage
\fancyhead[LH]{\zihao{-5}{\songti 重庆大学本科学生毕业论文（设计）}}
\fancyhead[RH]{\zihao{-5}{\songti 参考文献}}

\addcontentsline{toc}{section}{参考文献}
\renewcommand\refname{参考文献}

\begin{thebibliography}{99}
\setlength{\itemsep}{0pt}
\bibitem{unet} Ronneberger O, Fischer P, Brox T. U-Net: Convolutional Networks for Biomedical Image Segmentation[C]. MICCAI, 2015: 234-241.
\bibitem{unet3d} Cicek O, Abdulkadir A, Lienkamp S S, et al. 3D U-Net: Learning Dense Volumetric Segmentation from Sparse Annotation[C]. MICCAI, 2016: 424-432.
\bibitem{attention_unet} Oktay O, Schlemper J, Le Folgoc L, et al. Attention U-Net: Learning Where to Look for the Pancreas[J]. arXiv preprint arXiv:1804.03999, 2018.
\bibitem{nnunet} Isensee F, Jaeger P F, Kohl S A A, et al. nnU-Net: A Self-configuring Method for Deep Learning-based Biomedical Image Segmentation[J]. Nature Methods, 2021, 18(2): 203-211.
\bibitem{kamnitsas2017} Kamnitsas K, Ledig C, Newcombe V F J, et al. Efficient Multi-scale 3D CNN with Fully Connected CRF for Accurate Brain Lesion Segmentation[J]. Medical Image Analysis, 2017, 36: 61-78.
\bibitem{chen2021} Chen M, Geng C, Wang D, et al. Deep Learning-based Segmentation of Cerebral Aneurysms in 3D TOF-MRA using Coarse-to-Fine Framework[J]. arXiv preprint arXiv:2110.13432, 2021.
\bibitem{ou2022} Ou C, Qian Y, Chong W, et al. A Deep Learning-based Automatic System for Intracranial Aneurysms Diagnosis on Three-dimensional DSA Images[J]. Medical Physics, 2022, 49(11): 7038-7053.
\bibitem{strokenet} Irfan M, Malik K M, Ahmad J, et al. StrokeNet: An Automated Approach for Segmentation and Rupture Risk Prediction of Intracranial Aneurysm[J]. Computerized Medical Imaging and Graphics, 2023, 108: 102271.
\bibitem{xiong2025} Xiong Y, Ye M, Wang P, et al. Rupture Risk Prediction of Intracranial Aneurysm by Using Gene Expression Data Mining and Intelligent Optimization Algorithm[J]. International Journal of Bio-Inspired Computation, 2025, 26(1): 24-34.
\end{thebibliography}


\newpage
\fancyhead[LH]{\zihao{-5}{\songti 重庆大学本科学生毕业论文（设计）}}
\fancyhead[RH]{\zihao{-5}{\songti 附录A：关键配置与复现实验步骤}}

\addcontentsline{toc}{section}{附录A：关键配置与复现实验步骤}
\section*{附录A：关键配置与复现实验步骤}
本文实验通过YAML统一配置驱动，关键字段包括模型结构、补丁大小、优化器参数、风险任务特征筛选策略等。该方式便于复现实验并支持快速切换训练场景。

\subsection*{主要程序模块说明}
\begin{enumerate}
	\item 分割训练脚本：负责模型构建、损失计算、验证与检查点保存；
	\item 分割评估脚本：负责顺序补丁推理、体数据重建与NIfTI结果写出；
	\item 风险训练脚本：负责表格数据加载、交叉注意力模型训练、阈值搜索和测试评估；
	\item 数据模块：实现病例扫描、标签匹配、重采样、补丁采样与特征预处理。
\end{enumerate}

\subsection*{形态学特征说明}
系统可从分割掩膜中提取体积、表面积、长宽深、纵横比、球形度、紧致度等特征，作为风险预测候选输入。相关特征计算模块已在项目中实现，便于后续与临床统计分析结合。

\subsection*{项目复现实验步骤}
\begin{enumerate}
	\item 根据环境配置安装依赖并准备CTA影像与标签数据；
	\item 配置训练参数，运行分割训练脚本生成最优模型；
	\item 基于最优分割模型导出预测结果并准备风险表格数据；
	\item 运行风险训练脚本，得到测试集评估指标与报告文件。
\end{enumerate}


\newpage
\fancyhead[LH]{\zihao{-5}{\songti 重庆大学本科学生毕业论文（设计）}}
\fancyhead[RH]{\zihao{-5}{\songti 致谢}}

\addcontentsline{toc}{section}{致谢}
\section*{致\quad 谢}
本科毕业设计期间，首先要衷心感谢指导教师文静老师在课题方向、研究方法和论文撰写方面给予的耐心指导。老师严谨务实的科研态度和对细节的高标准要求，使我在项目实现和学术表达上都获得了显著提升。

感谢实验室同学在数据准备、环境配置和问题讨论中的帮助。每一次代码调试和实验复盘，都离不开大家的建议与支持。

感谢学院提供的计算资源与学习平台，使本课题能够顺利完成从模型训练到系统验证的全流程实践。

最后，感谢家人一直以来的理解与鼓励。你们的支持是我坚持完成本科阶段学习与毕业设计的动力来源。


\newpage
\thispagestyle{empty}

\addcontentsline{toc}{section}{原创性声明和使用授权书}
\begin{center}
\heiti \zihao{3}
原创性声明
\end{center}

\songti\zihao{-4}
郑重声明：所呈交的论文（设计）\underline{《  \hspace{6em}》}，是本人在导师的指导下，独立进行研究取得的成果。除论文（设计）中已经标注引用的内容外，本论文（设计）不包含其他人或集体已经发表或撰写过的作品成果。对本文的研究做出贡献的个人和集体，均已在文中以明确方式标明。本人完全意识到本声明的法律后果，并承诺因本声明而产生的法律结果由本人承担。

~\\
\begin{flushleft}
\begin{tabular}{l}
\songti\zihao{-4}
论文（设计）作者签名: \underline{\hspace{6em}}\\
\songti\zihao{-4}
日期：\underline{\hspace{6em}}
\end{tabular}
\end{flushleft}

~\\
\begin{center}
\heiti \zihao{3}
使用授权书
\end{center}

\songti\zihao{-4}
本论文（设计）作者完全了解学校有关保留、使用论文（设计）的规定，同意学校保留并向国家有关部门或机构送交论文（设计）复印件和电子版，允许论文（设计）被查阅和借阅。本人授权重庆大学将本论文（设计）的全部或部分内容编入有关数据库进行检索，可以采用影印、缩印或扫描等复制方式保存和汇编本论文（设计）。

~\\
\songti\zihao{-4}
本论文（设计）属于：\par
保\quad 密 $\Box$  \quad 在\underline{\qquad}年解密后适用本授权书\par
不保密 $\Box$

~\\
~\\
\begin{flushleft}
\songti\zihao{-4}
\begin{tabular}{l l}
论文（设计）作者签名：\underline{\hspace{6em}} \hspace{300mm}&指导教师签名：\underline{\hspace{6em}} \\
日期：\underline{\hspace{6em}} &日期：\underline{\hspace{6em}}\\
\end{tabular}
\end{flushleft}

\end{document} 